\section{System Architecture}
\label{sec:arch}

\begin{figure}[t]
\centering
\begin{tikzpicture}[
  font=\scriptsize,
  box/.style={draw, rounded corners=1.5pt, align=center, minimum height=0.55cm, inner sep=2pt, fill=white},
  store/.style={draw, cylinder, shape border rotate=90, aspect=0.18, align=center, minimum height=0.62cm, inner sep=1.5pt, fill=white},
  arr/.style={-{Stealth[length=1.6mm]}, semithick},
  darr/.style={{Stealth[length=1.6mm]}-{Stealth[length=1.6mm]}, semithick},
  node distance=0.28cm and 0.3cm]

% Row 1: user-facing
\node[box, minimum width=2.0cm] (ui) {Desktop UI\\(React/Tauri)};
\node[box, right=of ui, minimum width=2.1cm] (api) {Task API\\(FastAPI, loopback)};
\node[box, right=of api, minimum width=1.7cm] (obs) {Observatory\\(separate process)};

% Row 2: runner + agent
\node[box, below=0.45cm of api, minimum width=5.6cm, minimum height=0.95cm] (agent) {};
\node[anchor=north, font=\scriptsize\bfseries] at (agent.north) {LangGraph agent (Fig.~\ref{fig:graph})};
\node[box, font=\tiny, anchor=south west, minimum width=1.25cm] (plan) at ($(agent.south west)+(0.08,0.07)$) {plan / ground};
\node[box, font=\tiny, right=0.08cm of plan, minimum width=0.95cm] (exec) {execute};
\node[box, font=\tiny, right=0.08cm of exec, minimum width=0.95cm] (ver) {verify $\Phi$};
\node[box, font=\tiny, right=0.08cm of ver, minimum width=1.05cm] (rec) {recover};

\node[box, left=0.25cm of agent, minimum width=1.2cm] (runner) {Task runner\\checkpoints};

% Row 3: services
\node[box, below=0.5cm of agent.south west, anchor=north west, minimum width=1.7cm] (router) {Model router};
\node[box, right=0.2cm of router, minimum width=1.75cm] (sec) {Sanitizer,\\approval gate};
\node[box, right=0.2cm of sec, minimum width=1.7cm] (pool) {Browser pool\\(Playwright)};

% Row 4: local resources
\node[box, below=0.35cm of router, minimum width=1.7cm, fill=gray!8] (ollama) {Ollama (local)\\text + vision};
\node[store, below=0.35cm of sec] (db) {SQLite\\tasks, events};
\node[store, right=0.12cm of db] (ev) {Evidence\\PNG, JSON};
\node[store, left=0.12cm of db] (chroma) {Chroma};
\node[box, below=0.35cm of pool, minimum width=1.7cm, fill=red!6, xshift=0.45cm] (web) {Web pages\\(untrusted)};

% Edges
\draw[darr] (ui) -- (api);
\draw[arr] (api) -- (agent.north -| api);
\draw[arr] (api.south west) ++(0,0) -- (runner.north);
\draw[arr] (runner) -- (agent);
\draw[darr] (router.north) -- (router.north |- agent.south);
\draw[darr] (router) -- (ollama);
\draw[arr] (sec.north) -- (sec.north |- agent.south);
\draw[darr] (pool.north) -- (pool.north |- agent.south);
\draw[darr] (pool) -- (web);
\draw[arr] (agent.south) ++(0.3,0) -- ++(0,-0.25) -| (ev.north);
\draw[arr, dashed] (db.east) ++(0,0.1) -- ++(0.12,0) |- (obs.south);
\draw[arr, dashed] (api.east) -- (obs.west);

% Trust boundary: everything except the web runs on the user's machine
\begin{scope}[on background layer]
\node[draw, dashed, rounded corners=3pt, inner sep=3pt, fit=(ui)(obs)(runner)(ollama)(chroma)(ev)(pool), label={[font=\tiny, anchor=south east, inner sep=1pt]south east:user's machine}] {};
\end{scope}
\end{tikzpicture}
\caption{Components of Agentic Pilot. Everything inside the dashed boundary runs on the user's machine; model inference goes to a local Ollama server. Web pages are the main untrusted input. Dashed arrows are the telemetry feed read by the Observatory.}
\label{fig:arch}
\end{figure}


Fig.~\ref{fig:arch} shows the components and the trust boundary. A desktop client (React inside a Tauri shell) submits a natural-language task to a FastAPI service that binds only to the loopback interface. The task runner queues the task, starts one execution of the LangGraph agent, writes a checkpoint after every node so that a paused or interrupted task can resume, and enforces a wall-clock limit (5 minutes by default).

\paragraph{Agent graph} The agent is a compiled LangGraph state machine with eleven nodes (Fig.~\ref{fig:graph}). A typed state record carries the parsed intent, the task plan and step index, the current page manifest, the planned action, the action history, retry and model-call counters, routing decisions, and the final status. Nodes return partial updates to this record; conditional routers read only the status and a few flags. The graph has a single terminal node, \texttt{complete}, which persists the outcome, writes memory records and decides whether the browser context is kept open for the user.

\begin{figure}[t]
\centering
\begin{tikzpicture}[
  font=\scriptsize,
  n/.style={draw, rounded corners=2pt, minimum height=0.5cm, minimum width=1.25cm, inner sep=2pt, fill=white},
  g/.style={n, fill=gray!12, very thick},
  t/.style={n, fill=black!4, rounded corners=6pt},
  arr/.style={-{Stealth[length=1.5mm]}, semithick},
  lab/.style={font=\tiny, inner sep=1pt, fill=white},
  node distance=0.36cm and 0.42cm]

\node[n] (parse) {parse\_intent};
\node[n, right=of parse] (risk) {risk\_check};
\node[n, right=of risk] (auth) {auth\_check};
\node[n, below=of auth] (nav) {navigate};
\node[n, below=of nav] (dom) {extract\_dom};
\node[n, left=of dom] (ctx) {retrieve\_ctx};
\node[n, left=of ctx] (plan) {plan\_action};
\node[n, below=of plan] (exec) {execute};
\node[g, right=of exec] (ver) {verify ($\Phi$)};
\node[n, right=of ver] (recv) {recovery};
\node[t, below=0.45cm of ver] (done) {complete};

\draw[arr] (parse) -- (risk);
\draw[arr] (risk) -- (auth);
\draw[arr] (auth) -- (nav);
\draw[arr] (nav) -- node[lab, right] {ok} (dom);
\draw[arr] (dom) -- (ctx);
\draw[arr] (ctx) -- (plan);
\draw[arr] (plan) -- (exec);
\draw[arr] (exec) -- (ver);
\draw[arr] (ver) -- node[lab, above] {failed} (recv);
\draw[arr] (ver) -- node[lab, right] {done / blocked / approval} (done);
\draw[arr] (ver.north) -- node[lab, left, pos=0.35] {running} (ver.north |- dom.south) ;
\draw[arr] (recv.north) -- node[lab, right, pos=0.4] {retry} (recv.north |- dom.south);
\draw[arr] (recv.south) |- node[lab, pos=0.25, right] {exhausted} (done.east);
\draw[arr] (nav.east) -- ++(0.35,0) |- node[lab, pos=0.2, right] {fail} (recv.east);
\draw[arr] (parse.south) -- ++(0,-0.2) -| node[lab, pos=0.1, below] {fail} ($(recv.north)+(0.35,0)$);
\draw[arr] (risk.south) |- node[lab, pos=0.15, left] {needs approval} ($(plan.west)+(-0.25,0.35)$) |- (done.west);
\draw[arr, dashed] (dom.east) -- ++(0.55,0) |- node[lab, pos=0.9, above] {CAPTCHA} ($(done.east)+(0,0.1)$);
\end{tikzpicture}
\caption{Execution graph as compiled in \texttt{graph.py}. Only the shaded \texttt{verify} node may route to \texttt{complete} with status \textsc{Completed}; it does so when the completion predicate $\Phi$ holds. A detected CAPTCHA routes to \texttt{complete} with status \textsc{Blocked}.}
\label{fig:graph}
\end{figure}


\paragraph{Model layer} All language and vision calls go to an Ollama server at a configurable address, by default \texttt{127.0.0.1:11434}. A router chooses a model per call according to the role the call plays (Section~\ref{sec:routing}). A gateway enforces JSON output, validates it against Pydantic schemas and re-prompts on malformed responses.

\paragraph{Browser layer} A pool of Playwright contexts over one Chromium instance gives each task its own context and page. The executor exposes navigation, clicking, typing, key presses, scrolling, content extraction and screenshots, and reads back typed values after text entry.

\paragraph{Verification, recovery and evidence} A verification manager implements the completion predicate of Section~\ref{sec:formal}, and a recovery engine classifies failures and bounds retries (Section~\ref{sec:recovery}). An evidence manager writes, per task, before/after screenshots of every executed action, the DOM manifest, the verification and recovery records, a node-level trace, and a structured execution record per action. Task events are also stored in SQLite and streamed to observers.

\paragraph{Memory and security} Task outcomes are written to SQLite and to a local Chroma vector store (Section~\ref{sec:memory}). A sanitizer filters element text before it reaches the planner prompt, a risk gate pauses tasks that need human approval, a CAPTCHA detector halts execution, and a privacy auditor logs every model call with its destination and payload sizes (Section~\ref{sec:security}).

\paragraph{Observatory} A separate process subscribes to the agent's event stream over WebSocket (falling back to server-sent events) and reads the task database and evidence files to display plans, routing decisions, verification outcomes and screenshots. It issues only read queries and cannot send commands to the agent; the database connection is not opened in read-only mode, so this is a property of the code rather than an enforced constraint.

\paragraph{Trust boundaries} Two inputs cross into the agent from outside: the user's task text, which is trusted, and page content, which is not. Page content reaches the planner as element text, the page title and URL, and screenshots for the vision model. Outbound traffic consists of the browser's own requests to the sites the task visits; model traffic stays on the local host unless the user reconfigures the Ollama address.
